\section{Gauge covariance and the residual-energy identity}
\label{app:gauge}

Let $\mathcal T_{yx}$ denote the set of group-valued transports from $y$ to
$x$.  It is a left--right $\SU$ bitorsor.  Choosing frames identifies
$\mathcal T_{yx}$ with $\SU$, but a gauge transformation changes that
identification by
\[
  T\longmapsto g_yTg_x^{-1}.
\]
Because the metric is bi-invariant,
\[
 d(g_yTg_x^{-1},g_yT'g_x^{-1})=d(T,T').
\]
Hence the Fr\'echet functional and its infimum are intrinsic to the bitorsor.

For an incident plaquette oriented as $W_i=U_\ell S_i$,
\begin{align*}
 d(I,U_\ell S_i)
 &=d(U_\ell^{-1},S_i),
\end{align*}
where left multiplication by $U_\ell^{-1}$ was applied to both arguments.
Summing over $i$ proves Theorem~\ref{thm:residual}.  Compactness of $\SU$
guarantees that the infimum is attained.

For the diameter witness, if $D\le\eta^2$ and $C_*$ is a minimizer, then each
nonnegative summand satisfies $d(C_*,S_i)^2\le D\le\eta^2$.  Thus all staples
belong to $\overline B(C_*,\eta)$.  Taking the contrapositive proves the stated
form.
